% macos-security-architecture.tex — the macOS protection layers from boot to runtime on Apple
% silicon (the only Mac platform from macOS 27): Secure Boot + LocalPolicy (Full / Reduced /
% Permissive, 1TR), SIP, the Signed System Volume, the hardware runtime protections (KIP, SPTM,
% MIE), Gatekeeper + quarantine + notarization at first open, XProtect and the macOS 26.4
% terminal / script protections, Background Security Improvements, TCC, hardened runtime vs App
% Sandbox, FileVault, Lockdown Mode.
% Pictures: what happens when a downloaded app is opened; the Apple-silicon boot chain.
% Sources (checked 2026-10-09):
%   support.apple.com/guide/security — Apple Platform Security (covers the 26.5 releases; revised Aug 2026):
%     boot-process-secac71d5623 (Boot ROM → LLB → iBoot; LLB checks the LocalPolicy anti-replay value;
%       Full / Reduced / Permissive; any signature failure → recoveryOS; AuxKC + SCIP; iBoot checks the
%       SSV root hash)
%     startup-disk-security-policy-control-sec7d92dc49f (Full = personalised ECID signature, security
%       epochs; Reduced = global signature, required for kexts, AuxKC hash in LocalPolicy, restart;
%       Permissive = locally SE-signed boot objects, required to turn SIP off, CLI only; SIP store moved
%       from NVRAM to LocalPolicy; power-button recoveryOS; no firmware password)
%     contents-a-localpolicy-file-mac-apple-silicon-secc745a0845 (1TR = single press-and-hold; OIK =
%       "Ownership"; nsih SHA-384; auxp/auxi/auxr; coih via kmutil configure-boot; sip0–sip3; which
%       fields are mutable in 1TR only)
%     system-integrity-protection-secb7ea06b49 (MAC on every process; KIP separate; OS X 10.11)
%     signed-system-volume-security-secd698747c9 (SHA-256 in the metadata tree, seal, read-path check,
%       boot halts on failure, APFS snapshots, no SSV-off while FileVault is on)
%     operating-system-integrity-sec8b776536b (KIP, SCIP, PAC; SPTM + TXM on M2 or later replace PPL;
%       neither forces signed-only code on macOS; MIE = EMTE synchronous on A19 / M5 or later)
%     gatekeeper-and-runtime-protection-sec5599b66df (Developer ID + notarized by default; first-open
%       malware check regardless of origin; override unless MDM restricts; randomized read-only locations)
%     protecting-against-malware-sec469d47bd8 (revocation tickets via CloudKit, more frequent than XProtect
%       updates; XProtect YARA, daily check, 10.15 triggers, Trash; remediation + behavioural engines;
%       Endpoint Security events for Gatekeeper bypass + XProtect detections in macOS 15)
%     terminal-and-script-protections-sece3b202c4b (macOS 26.4: paste warning + its conditions, 24 h grace,
%       pasteboard command blocking, AppleScript / JXA scanning, in-memory scripts not overridable)
%     background-security-improvements-sec87fc038c2 (26.1; cryptexes on the preboot volume; no reseal;
%       Safari relaunch; patch + antipatch)
%     controlling-app-access-to-files-secddd1d86a6 (10.15 folders; full-disk access since 10.13)
%     volume-encryption-with-filevault-sec4c6dc1b6e (AES-XTS; keys only in the Secure Enclave; UID-only
%       key while off; enabling is immediate; 24-character recovery key; iCloud Keychain sync)
%     data-protection-overview-secf6276da8a (macOS 26.4: FileVault on by default; Class C volume key)
%     securely-extending-the-kernel-sec8e454101b (kexts: 1TR → Reduced; system extensions in user space)
%     lockdown-mode-security-sec2437264f0 (macOS 13, more in 14)
%   support.apple.com/en-us/102149 (SIP paths; startup disk) · /102445 (Open Anyway flow, alert
%     meanings) · /102657 + /126604 (BSI from 26.1; first one macOS 26.3.1 (a), 17 Mar 2026; built-in
%     disk only) · /105120 (Lockdown Mode: profiles, accessories) · /127255 (macOS 27: Apple silicon
%     only) · /149035 (macOS 27 security content: released 14 Sep 2026) · /124963 (Tahoe 26: FileVault unlock over SSH; 26.2 MDM grants shown) · /148830 (macOS 27:
%     MDM launch control, consolidated consent, PPPC Accessibility notice, tccutil list, no direct
%     TCC database access)
%   support.apple.com/guide/deployment/privacy-preferences-policy-control-payload-dep38df53c2a (camera,
%     microphone, screen capture: deny only; most restrictive payload wins)
%   developer.apple.com/news/?id=saqachfa (macOS 15: no Control-click override) · /news/?id=w4atic4c
%     (original Developer ID Sub-CA expires 1 Feb 2027; .pkg stops installing; G2)
%   developer.apple.com/documentation: security/hardened-runtime · security/app-sandbox ·
%     security/accessing-files-from-the-macos-app-sandbox (container, Powerbox panels, bookmarks, macOS 14
%     container consent, no entitlement for Full Disk Access, EPERM vs EACCES) ·
%     security/notarizing-macos-software-before-distribution (plug-ins 10.15) ·
%     security/resolving-common-notarization-issues (get-task-allow) · entitlements com.apple.security.cs.*
%     (library validation = Apple or same Team ID; MAP_JIT) · technotes/tn3127 (designated requirement =
%     the identity TCC records) · technotes/tn3179 (local network, macOS 15) ·
%     macos-release-notes/macos-big-sur-11_0_1-universal-apps-release-notes (arm64 must be signed)
%   developer.apple.com/videos/play/wwdc2022/10096 (Ventura: all notarized apps integrity-checked;
%     NSUpdateSecurityPolicy, App Management)
%   developer.apple.com/forums/thread/706442 + /130560 (Apple DTS: curl / scp do not quarantine,
%     tar / unzip do not propagate, Archive Utility does; syspolicy_check since macOS 14)
% Not repeated: the code-signature format (code-signing-internals); notarytool + packaging
% (notarization-distribution); keychain (keychain-secure-enclave); launchd + SMAppService
% (macos/xnu-launchd-xpc); the iOS sandbox + PAC detail (os/ios-sandbox-security-model).
% @labels: area=macos kind=concept platform=apple topic=security,platform-apis discussion=323
% @tags: secure-boot, localpolicy, sip, signed-system-volume, gatekeeper, quarantine, notarization, developer-id, xprotect, tcc, hardened-runtime, app-sandbox, entitlements, filevault, lockdown-mode, endpoint-security
\documentclass[8pt]{extarticle}
\usepackage{printup-sheet}
\usepackage{array}

\lstdefinelanguage{ShSheet}{
  morekeywords={xattr,spctl,syspolicy_check,codesign,csrutil,tccutil,fdesetup},
  sensitive=true, morecomment=[l]{\#},
  literate={--}{{\hbox{-}\hbox{-}}}2 {->}{{\hbox{-}\hbox{>}}}2}
\lstset{basicstyle=\ttfamily\scriptsize, aboveskip=2pt, belowskip=2pt}

\newcommand\ct[1]{\texttt{#1}}
\newcommand\hd[1]{\par\vspace{2pt}\noindent{\bfseries\color{sheetBlue}#1}\par\vspace{1pt}}

\tikzset{
  lbl/.style={font=\tiny, text=black!75, inner sep=1pt, align=center},
  st/.style={box, font=\tiny, text width=20mm, minimum height=15mm, align=center, inner sep=1.5pt},
  bc/.style={box, font=\tiny, text width=15.5mm, minimum height=10mm, align=center, inner sep=1.5pt},
  bad/.style={box, font=\tiny, draw=sheetRed, fill=sheetRed!7, align=center, inner sep=1.5pt},
  yes/.style={->, thick, draw=sheetGreen!70!black},
  no/.style={->, thick, draw=sheetRed},
}

\begin{document}

\sheettitle{macOS security architecture — boot, seal, first launch, runtime}{macos · folio}

\oneliner{Each layer catches what the one before cannot: \textbf{Secure Boot} starts only an
Apple-signed macOS, as its \textbf{LocalPolicy} allows; \textbf{SIP} keeps the \textbf{sealed}
system volume (\textbf{SSV}) read-only even for \ct{root}; a \textbf{quarantined} download meets
\textbf{Gatekeeper} (signature, Developer ID, \textbf{notarization}, \textbf{XProtect}) at first
open; then the \textbf{hardened runtime} blocks injection, the \textbf{App Sandbox} confines,
\textbf{TCC} asks per resource, \textbf{FileVault} encrypts. macOS 27 (Sep 2026): Apple silicon only.}

\vspace{1pt}
\noindent\begin{tikzpicture}[sheet]
  \node[font=\bfseries\scriptsize, anchor=west, text=sheetBlue] at (-0.05,1.0) {First open of a downloaded app — what Gatekeeper checks (Apple does not document the order)};
  \node[st, fill=sheetGreen!10, draw=sheetGreen!70!black] (s1) at (1.05,0) {\textbf{1 download}\\Safari, AirDrop set \ct{com.apple.quarantine}; Archive Utility keeps it\\\textcolor{sheetRed}{\ct{curl}, \ct{scp}: none; \ct{tar}, \ct{unzip} drop it}};
  \node[st] (s2) at (3.4,0) {\textbf{2 open}\\quarantined $\to$ Gatekeeper; may run it from a \textbf{randomized read-only} path, so plug-ins beside it don't load};
  \node[st] (s3) at (5.75,0) {\textbf{3 signature valid?}\\code + sealed resources unchanged since signing};
  \node[st] (s4) at (8.1,0) {\textbf{4 identified?}\\Developer ID or App Store signed; certificate not revoked};
  \node[st] (s5) at (10.45,0) {\textbf{5 notarized?}\\ticket stapled or found online; no revocation ticket};
  \node[st] (s6) at (12.8,0) {\textbf{6 known malware?}\\XProtect YARA rules, on every first open, whatever the origin};
  \node[st, fill=sheetGreen!10, draw=sheetGreen!70!black] (s7) at (15.15,0) {\textbf{7 consent}\\``downloaded from the Internet — open?'' asked once};
  \foreach \a/\b in {s1/s2,s2/s3,s3/s4,s4/s5,s5/s6,s6/s7} \draw[yes] (\a) -- (\b);
  \node[bad, text width=19mm] (d3) at (5.75,-1.45) {``can't be opened'': modified or damaged};
  \node[bad, text width=42mm] (d45) at (9.275,-1.45) {\textbf{blocked}: developer cannot be verified / Apple cannot check it. \textbf{macOS 15}: no Control-click; Privacy \& Security $\to$ \textbf{Open Anyway} $\to$ Open = saved exception};
  \node[bad, text width=19mm] (d6) at (12.8,-1.45) {``will damage your computer''; known malware to the Trash};
  \draw[no] (s3) -- node[lbl, right]{no} (d3);
  \draw[no] (s4) -- node[lbl, right]{no} (s4 |- d45.north);
  \draw[no] (s5) -- node[lbl, right]{no} (s5 |- d45.north);
  \draw[no] (s6) -- node[lbl, right]{hit} (d6);
  \node[lbl, anchor=north west, align=left, text width=38mm] at (-0.05,-0.86) {\textbf{macOS 13+}: every notarized app is integrity-checked, quarantined or not; \textbf{App Management}: only the same team (or an \ct{NSUpdateSecurityPolicy}) may modify an app without the user's OK.};
  \node[lbl, anchor=north west, align=left, text width=19mm] at (14.1,-0.86) {replay the verdict:\\\ct{syspolicy\_check}\\\ct{distribution} (14+)};
  \fill[sheetBlue!12] (-0.05,-2.42) rectangle (16.55,-2.08);
  \node[font=\tiny, anchor=west, text width=164mm] at (0,-2.25) {\textbf{every launch}: arm64 code needs a valid signature (ad hoc suffices) · hardened runtime · sandbox if entitled · \textbf{TCC} at first camera / mic / Documents use · XProtect rescans changed apps};
\end{tikzpicture}

\begin{multicols}{2}
\raggedright
\footnotesize\setstretch{1.0}

\hd{Secure Boot on Apple silicon}
\noindent\begin{tikzpicture}[sheet]
  \node[bc, fill=black!6, draw=black!60] (rom) at (0.8,0) {\textbf{Boot ROM}\\in the chip: first link of trust};
  \node[bc] (llb) at (2.78,0) {\textbf{LLB}\\loads LocalPolicy, checks its SE signature + anti-replay value};
  \node[bc] (ib) at (4.76,0) {\textbf{iBoot}\\checks AuxKC + SSV root hash; SCIP locks the kernel collections};
  \node[bc, fill=sheetGreen!10, draw=sheetGreen!70!black] (k) at (6.74,0) {\textbf{XNU}\\mounts the sealed volume; KIP locks kernel code};
  \draw[yes] (rom) -- (llb); \draw[yes] (llb) -- (ib); \draw[yes] (ib) -- (k);
  \node[bad, text width=76mm] (rec) at (3.77,-1.05) {any signature failure $\to$ \textbf{recoveryOS} to repair · seal mismatch $\to$ boot halts, reinstall macOS};
  \draw[no] (llb) -- (llb |- rec.north); \draw[no] (ib) -- (ib |- rec.north);
\end{tikzpicture}
{\scriptsize\setlength\tabcolsep{2pt}
\begin{tabular}{@{}>{\raggedright\arraybackslash}p{11mm}>{\raggedright\arraybackslash}p{66mm}@{}}
\toprule
\textbf{Full} & default, as on iOS: install fetches a \textbf{personalised} signature (ECID) from Apple's signing server, so only the newest macOS at install time boots; the server refuses old security epochs \\
\textbf{Reduced} & trusts the \textbf{global} signature shipped with macOS (older versions); required for third-party \textbf{kexts} — merged into an \textbf{AuxKC} whose hash is in the LocalPolicy, loaded after a restart, never on demand \\
\textbf{Permissive} & chain still verified, but iBoot also accepts objects the Secure Enclave signed locally (a custom XNU kernel collection); required to turn \textbf{SIP off}; command line only \\
\bottomrule
\end{tabular}}\par
\textbf{LocalPolicy}: an Image4 file per installed macOS, signed by the \textbf{Owner Identity Key}
in the Secure Enclave (``Ownership''). Full $\leftrightarrow$ Reduced in Startup Security Utility.
\textbf{1TR} (paired One True recoveryOS) = power button single press-and-hold; NVRAM, a double
press or a boot error reach plain recoveryOS, which may set Reduced but not the SIP bits or a
custom kernel. Malware cannot hold a button, so these Macs neither need nor support a
\textbf{firmware password}. Fields: \ct{nsih} SHA-384 of the macOS manifest (iBoot, trust cache, kernel
collection, SSV root hash) · \ct{auxp}/\ct{auxi} approved kexts + AuxKC · \ct{coih} custom
kernel (\ct{kmutil configure-boot}) · \ct{sip0} SIP bits (were NVRAM) · \ct{sip1} accept a bad
SSV root hash · \ct{sip2} leave kernel memory writable · \ct{sip3} drop the boot-args allow-list.
Prefer \textbf{system extensions} (10.15, user space: DriverKit, Endpoint Security, Network
Extension) to kexts.

\hd{SIP — System Integrity Protection (OS X 10.11)}
\begin{itemize}
  \item Mandatory access control (same kernel mechanism as the sandbox) on \emph{every}
        process, sandboxed or \ct{root}: \ct{/System}, \ct{/usr}, \ct{/bin}, \ct{/sbin},
        \ct{/var} and preinstalled apps are writable only by Apple-signed processes with special
        entitlements (software update, installers). Third parties write \ct{/Applications},
        \ct{/Library}, \ct{/usr/local}. Also stops software choosing the startup disk.
  \item Denial is \ct{EPERM}; POSIX permissions and ACLs fail with \ct{EACCES}.
  \item \ct{csrutil disable} only in 1TR, by a user holding Ownership $\to$ Permissive; kept
        in the LocalPolicy, not NVRAM. SIP off also stops kext signature checks when the AuxKC is
        built. Xcode debugging never needs it off; dev DriverKit installs may.
\end{itemize}

\hd{Signed System Volume (macOS 11)}
Each file's SHA-256 sits in the file-system metadata tree, itself hashed up to one root — the
\textbf{seal}, signed by Apple and re-measured on the Mac at install and update. Data is hashed on every
read; a mismatch is never returned. iBoot verifies the seal before the kernel starts. Built on
APFS snapshots: a failed update returns to the old system without reinstalling. The system volume
needs no encryption any more, and SSV cannot be turned off while FileVault is on. \textbf{SIP stops
writes; SSV detects them.}

\hd{Hardware runtime protections}
{\scriptsize\setlength\tabcolsep{2pt}
\begin{tabular}{@{}>{\raggedright\arraybackslash}p{13mm}>{\raggedright\arraybackslash}p{64mm}@{}}
\toprule
\textbf{KIP} & after kernel start-up the memory controller refuses writes to kernel + kext code \\
\textbf{SCIP} & the same lock for coprocessor firmware (and the AuxKC region) \\
\textbf{PAC} & signed pointers against memory-corruption exploits \\
\textbf{SPTM + TXM} & M2 or later: SPTM runs above the kernel and guards page tables even against kernel write access; on macOS it does \emph{not} force signed-only code \\
\textbf{MIE} & M5 or later: Enhanced MTE (synchronous) + secure allocators, always on for the kernel and key attack surfaces \\
\bottomrule
\end{tabular}}

\hd{FileVault}
Apple silicon: internal storage is \emph{always} AES-XTS encrypted, keys handled only in the
Secure Enclave. Off: the key is guarded by the hardware UID alone, so turning it on is instant
(UID + password; old key retired). \textbf{On by default from macOS 26.4}
(opt-out). Recovery key: 24 characters, synced by iCloud Keychain if used. macOS 26: unlock over
SSH after a restart when Remote Login is on.

\hd{Lockdown Mode (macOS 13, more in 14)}
For targeted users: complex web tech off (per-site exceptions), no new profiles or MDM
enrolment; on Apple-silicon laptops an accessory needs unlock + approval.

\columnbreak

\hd{Gatekeeper \& notarization — the rules}
\begin{itemize}
  \item Default: outside the App Store = \textbf{Developer ID} signature + \textbf{notarized}
        (10.15+). Settings: App Store only / + identified developers; per-app override unless MDM
        forbids; Gatekeeper can be switched off entirely.
  \item Notarization = Apple's automated malware + signing scan, not App Review. Revocation
        tickets reach Macs by CloudKit, far more often than XProtect rules. Since 10.15 a
        quarantined plug-in loads only if notarized or user-approved.
  \item Arm64 code must be signed — ad hoc passes the kernel, never Gatekeeper; Rosetta x86 code
        is exempt (general-purpose Rosetta through macOS 27).
  \item macOS 15: Endpoint Security reports Gatekeeper overrides + XProtect hits. macOS 27: MDM
        decides which apps and binaries may run; App Attest on the Mac.
  \item 1 Feb 2027: the original Developer ID Sub-CA expires — \ct{.pkg}s signed under it stop
        installing (re-sign with G2); timestamped, notarized apps keep running.
\end{itemize}

\hd{XProtect \& the newer layers}
\begin{itemize}
  \item \textbf{XProtect}: YARA rules updated apart from OS updates (checked daily); scans at
        first launch, when an app changes, when rules update; a hit is blocked and moved to the
        Trash. Plus a periodic \textbf{remediation} engine and a \textbf{behavioural} one.
  \item \textbf{Terminal \& scripts (26.4)}: (1) Terminal asks before running text pasted from a
        browser, Messages, Mail… if unused 30+ days and no Xcode / IDE is found (24 h grace after
        setup). (2) Paste into \emph{any} terminal: XProtect follows the process tree, checks its
        hosts against Safe Browsing, blocks bad ones. (3) AppleScript + JXA scanned in OpenScripting;
        in-memory scripts cannot be overridden.
  \item \textbf{Background Security Improvements} (26.1+): Safari / WebKit / library fixes as
        sealed \textbf{cryptexes} on the preboot volume — no SSV reseal; active after a restart
        (Safari: a relaunch). Patch + antipatch: removal returns to the base update. First one:
        macOS 26.3.1 (a), 17 Mar 2026. Built-in startup disk only.
\end{itemize}

\hd{TCC — Transparency, Consent \& Control}
\begin{itemize}
  \item Per app, per resource: camera, microphone, screen recording, Input Monitoring,
        Accessibility, Automation, \textbf{Files \& Folders} (Desktop, Documents, Downloads, iCloud
        Drive, network + removable volumes; 10.15). \textbf{Full Disk Access} (10.13): the user
        adds the app in System Settings; no entitlement grants it. Hardened apps also need the
        resource entitlement (\ct{device.camera}…) beside the usage string.
  \item Identity = the \textbf{designated requirement}: ad hoc = this build only (re-prompts per
        build); Apple Development $\neq$ Developer ID; Xcode-signed App Store and Developer ID
        builds share grants. macOS 14: another app's container needs consent.
  \item \textbf{Local Network} (macOS 15): per user; launchd daemons, \ct{root} and Terminal / SSH
        tools pass; agents don't; credited to the responsible app; no MDM control.
  \item MDM (PPPC): camera, microphone, screen capture can only be \emph{denied}; the most
        restrictive payload wins. macOS 27: no direct TCC-database reads; \ct{tccutil list}
        (managed); MDM-granted Accessibility is announced, user-revocable.
\end{itemize}

\hd{Hardened runtime vs App Sandbox {\normalfont(prefix \ct{com.apple.security.})}}
{\scriptsize\setlength\tabcolsep{2pt}
\begin{tabular}{@{}>{\raggedright\arraybackslash}p{9mm}>{\raggedright\arraybackslash}p{37mm}>{\raggedright\arraybackslash}p{31mm}@{}}
\toprule
 & \textbf{hardened runtime} & \textbf{App Sandbox} \\ \midrule
stops & JIT / W+X memory, \ct{DYLD\_*}, libraries not Apple- or team-signed, memory tampering & all outside its container: files, network, devices, other apps \\
needed for & \textbf{notarization} (no \ct{get-task-allow}) & \textbf{Mac App Store} \\
opt-outs & \ct{cs.allow-jit}, \ct{cs.disable-library-validation} (Gatekeeper checks harder) & \ct{files.user-selected.read-write}, \ct{network.client} \\
\bottomrule
\end{tabular}}\par
Container: \ct{\~{}/Library/Containers/<id>}. Open / save panels run \textbf{out of process} and
widen the sandbox to the pick; a \textbf{security-scoped bookmark} keeps it.

\hd{Inspect every layer}
\begin{lstlisting}[language=ShSheet]
xattr -p com.apple.quarantine Tool.app  # present: Gatekeeper runs
syspolicy_check distribution Tool.app   # 14+; older: spctl -a -t exec
codesign -d --entitlements - Tool.app   # runtime exceptions
csrutil status                          # SIP (change it only in 1TR)
fdesetup status                         # FileVault is On.
\end{lstlisting}

\hd{Traps}
\begin{itemize}
  \trap{\ct{xattr -d com.apple.quarantine} skips Gatekeeper for that copy — \textbf{notarize}
        instead. \ct{strip} / \ct{install\_name\_tool} after linking break the signature.}
  \trap{SIP $\neq$ sandbox $\neq$ TCC: SIP guards the OS from \ct{root}, the sandbox confines one
        app, TCC asks the user.}
\end{itemize}

\end{multicols}

\noindent{\footnotesize\color{sheetGrey}\raggedright\textit{Related:} Code signing internals · Notarization
\& distribution channels · \texttt{keychain-secure-enclave} · macOS keychains \& the security
CLI · \texttt{xnu-launchd-xpc} · \texttt{ios-sandbox-security-model} · \texttt{mach-o-and-dyld} ·
\texttt{macos-power-user-cli} · \texttt{linux-security-permissions}\par}

\end{document}
